\begin{figure}[htbp]
  \centering
  \begin{tikzpicture}[
    conditioned cell/.style={
      draw=black!45, dashed, fill=black!7,
      minimum width=25mm, minimum height=10mm,
      inner sep=3pt
    },
    operator cell/.style={
      draw=black, fill=white, line width=.6pt,
      minimum width=25mm, minimum height=10mm,
      inner sep=3pt
    },
    operator arrow/.style={-{Stealth[length=2.2mm]}, line width=.7pt},
    operator label/.style={fill=white, inner sep=2pt, font=\small}
  ]
    \node at (0,1) {$r$};
    \node at (4,1) {$s$};
    \node at (8,1) {$t$};

    \node[anchor=east] at (-1.8,0) {$P$};
    \node[anchor=east] at (-1.8,-2) {$A$};
    \node[anchor=east] at (-1.8,-4) {$B$};

    \node[conditioned cell] (Pr) at (0,0) {$z=Y_{P,r}$};
    \node[conditioned cell] (Ps) at (4,0) {$c=Y_{P,s}$};
    \node[conditioned cell] (Pt) at (8,0) {$d=Y_{P,t}$};
    \node[conditioned cell] (Ar) at (0,-2) {$x=Y_{A,r}$};
    \node[conditioned cell] (Br) at (0,-4) {$y=Y_{B,r}$};

    \node[operator cell] (As) at (4,-2) {$H_{A,s}$};
    \node[operator cell] (At) at (8,-2) {$H_{A,t}$};
    \node[operator cell] (Bs) at (4,-4) {$H_{B,s}$};
    \node[operator cell] (Bt) at (8,-4) {$H_{B,t}$};

    \draw[operator arrow] (At.west) --
      node[operator label,above] {$L_A(x)$} (As.east);
    \draw[operator arrow] (Bt.west) --
      node[operator label,above] {$L_B(y)$} (Bs.east);
    \draw[operator arrow] (Bs.north) --
      node[operator label,left] {$J_s(c)$} (As.south);
    \draw[operator arrow] (Bt.north) --
      node[operator label,right] {$J_t(d)$} (At.south);
  \end{tikzpicture}
  \caption{The grouped operator square at a fixed parameter $E$.
    The five shaded cells record the conditioned values $z,c,d,x,y$.
    Each arrow acts on functions in its source Hilbert space by
    conditional expectation in the indicated row or column law.
    Reversing an arrow represents its adjoint.}
  \label{fig:operator-square}
\end{figure}
